\subsection{CIC Filter}\label{subsec:CIC_Filter}

Ein CIC-Filter oder Cascaded-Integrator-Comb-Filter ist ein digitales Filter, welches nebst der Filterung eines Signals zusätzlich deren Abtastfrequenz verändert. Abbildung \ref{img:CIC_Filter} zeigt ein Dezimations-CIC-Filter. Dieses verkleinert die Abtastfrequenz am Ausgang um den Faktor \(R\). Die zweite Form ist ein Interpolation-CIC-Filter. Dieses vergrössert die Abtastfrequenz um den Faktor \(R\).

In Abbildung \ref{img:CIC_Filter} ist zu sehen, dass das Filter in drei Stufen unterteilt ist. Links ist ein Integrator-Filter zu sehen, welches einen anliegenden Wert mit einem verzögerten Wert aufaddiert oder anders gesagt, das Signal integriert. Anschliessend folgt ein Dezimierer, welcher das Signal um den Faktor R unterabtastet. Zuletzt folgt ein Combfilter. Dieses nimmt den aktuellen Wert und subtrahiert den alten Wert. Ein Interpolation-CIC-Filter erhält man durch tauschen von Integrator-Filtern mit Comb-Filtern und umgekehrt. Weiter ist nun in der Mitte eine Überabtastung nötig.\\
Es ist möglich, mehrere Integrator- und Combfilterpaare zusammen zu schalten, um eine grössere Dämpfung höherer Frequenzen zu bewirken. Das CIC-Filter benötigt jedoch, um zu funktionieren, eine gewisse Anzahl Bits in den Speichern der Verzögerungselemente, welche grösser als die Anzahl Eingangsbits ist. Dieser Effekt nennt sich Bit-Growth und die zusätzliche Anzahl Bits lässt sich wie folgt berechnen \cite{cic_a} \cite{cic_h}.

\begin{equation}
B_+ = \big \lceil Nlog_2RM \big \rceil
\label{equ:cic_bitgrowth}
\end{equation}

Dabei entspricht N der Ordnung des Filters oder wie viele Integrator-Comb-Filterpaare das Filter hat. R ist der zuvor genannte Dezimationsfaktor und M ist die Verzögerung der Speicherelemente.

Das CIC-Filter verstärkt das Eingangssignal zudem um einen Faktor \(G\). Dieser lässt sich wie folgt berechnen:

 \begin{equation}
 G = (R\cdot M)^N
 \label{equ:cic_gain}
 \end{equation}
 
Um nun den Faktor zu berechnen, um welchen man den Ausgang eines CIC-Filter multiplizieren muss, um den Zahlenbereich voll auszunutzen, kann folgende Formel eingesetzt werden:

\begin{equation}
G_+ = G/2^{B_+}
\label{equ:cic_gain+}
\end{equation}

Ein Problem, welches die CIC-Filter mit sich bringen, ist, dass sie in bestimmten Situationen Aliasing erzeugen. Dieses entsteht wenn sich Signalkomponenten zu nahe an den Nullstellen des Filters befinden \cite{CIC_Aliasing}. In Abbildung \ref{img:CIC_Filter_plot} sieht man sehr schön, dass sehr auf die Frequenz des Signals im Vergleich zum verwendeten CIC-Filter geachtet werden muss.


\begin{figure}[h]
	\centering
	\includegraphics[width=\textwidth]{CIC_Filter.pdf}
	\caption{Aufbau eines CIC-Filters N-ter Ordnung}
	\label{img:CIC_Filter}
\end{figure}

\begin{figure}[h]
	\centering
	\includegraphics[width=0.9\textwidth]{CIC_Filter_plot.pdf}
	\caption{Amplitudengang und Phasengang eines CIC-Filters (M = 1, R = 5; N = 2)}
	\label{img:CIC_Filter_plot}
\end{figure}

\clearpage
